\section{失败复盘框架：把事故转化为系统能力}

在真实部署中，失败不可避免。讲者给出的精神是“不断接触现实并修正系统”。这一节总结一个复盘框架，帮助团队把事故转化为可复用资产，而不是一次性救火。

\subsection{复盘对象与分级}

复盘首先要按风险而不是按技术模块分级，因为同一个模型错误可能只是措辞不佳，也可能触发不可逆业务动作。分类的目标是决定响应时限、遏制范围和批准层级，并让最严重事件优先进入策略与评测回归，而不是把所有失败塞进同一个 backlog。

建议把失败分为三类：
\begin{itemize}[leftmargin=1.5em]
    \item A 类：安全与合规事故（越权建议、敏感信息泄露）；
    \item B 类：业务动作事故（错误退款、漏执行、重复执行）；
    \item C 类：体验事故（高延迟、话轮错乱、重复追问）。
\end{itemize}
A 类必须 24 小时内完成根因与遏制；B 类 48 小时内完成补偿与修复；C 类纳入周度质量迭代。

\subsection{一页纸复盘模板}

完成分级后，团队需要用统一结构把事实、根因、防线和修复连起来。模板应足够短，让值班团队能在事故后立即填写；也要足够具体，避免“模型偶发错误”这种无法行动的结论。尤其要区分触发错误的近因与允许错误扩大影响的系统性缺口。

\begin{table}[H]
\centering
\caption{Agent 事故一页纸复盘模板}
\begin{tabularx}{\textwidth}{>{\raggedright\arraybackslash}p{3cm} >{\raggedright\arraybackslash}X}
\toprule
模块 & 应填写内容 \\
\midrule
事件摘要 & 发生时间、影响范围、用户影响、检测方式 \\
根因分析 & 模型行为、工具行为、策略行为分别发生了什么 \\
防线回顾 & 输入守卫、执行守卫、输出守卫哪一层失效 \\
补偿动作 & 对用户与业务系统采取了哪些纠偏动作 \\
长期修复 & 评测集新增项、规则新增项、流程新增项 \\
验证结果 & 修复后在仿真与灰度中的表现对比 \\
\bottomrule
\end{tabularx}
\end{table}

\begin{importantbox}{复盘的目标}
复盘不是追责会，而是“把一次失败转化为未来不再同类失败的系统升级”。只修 prompt 不修机制，事故会以新形态重现。
\end{importantbox}

\subsection{三类高频反模式}

\begin{enumerate}[leftmargin=1.8em]
    \item \textbf{反模式 1：} 把所有问题归因于“模型不够强”。
    实际上，很多事故由工具协议不稳、策略缺口、数据脏读引起。
    \item \textbf{反模式 2：} 出事后临时加规则但不进评测集。
    结果是下一版回归时旧问题复发，团队进入重复救火。
    \item \textbf{反模式 3：} 只修线上，不修开发流程。
    没有门禁和模板，事故知识无法沉淀，组织学习停滞。
\end{enumerate}

\begin{warningbox}{“热修复文化”会吞噬迭代速度}
如果团队长期依赖线上热修复而缺乏制度化复盘，系统会快速积累隐性技术债：规则互相冲突、行为不可解释、发布风险不断升高，最终拖慢全部创新节奏。
\end{warningbox}

\subsection{本章小结}

高质量 Agent 团队不是“从不失败”，而是“失败后系统能力单调上升”。复盘机制的成熟度，直接决定了产品在复杂环境下的长期进化速度。

